\section{总结与延伸}

本讲把自我改进从 inference-time selection 推进到 train-time distribution shift。STaR 从少量 rationale 出发，通过答案过滤和 rationalization 自举推理数据；DeepSeekMath 说明高质量领域数据与 practical RL 同样重要，并以 GRPO 用组内相对 reward 移除 critic；DAPO 则针对长 CoT RL 的真实病理，引入非对称 clipping、dynamic sampling、token-level loss 和 soft overlong punishment。

三个层次形成递进关系：

\begin{enumerate}
    \item \textbf{数据层}：生成并选择有学习价值的 reasoning trajectory；
    \item \textbf{目标层}：用相对 reward 提高正确轨迹概率；
    \item \textbf{系统层}：维持 entropy、有效梯度、长度与在线采样吞吐。
\end{enumerate}

下一讲转向“通过搜索实现自我改进”：AlphaCode 如何用海量采样、聚类和打分解决竞赛编程，Search-O1 如何在 reasoning 过程中按不确定性触发检索，并与 RL-based Search-R1 比较。

\end{document}
